% Generated draft: requires mathematical and source review.
\documentclass[11pt]{article}
\usepackage[margin=25mm]{geometry}
\usepackage{fontspec}
\setmainfont{Latin Modern Roman}
\usepackage{amsmath,amssymb,mathtools,unicode-math}
\setmathfont{Latin Modern Math}
\usepackage{enumitem,needspace,xurl,hyperref}
\setlength{\parindent}{0pt}
\setlength{\parskip}{5pt}
\setlength{\emergencystretch}{3em}
\newcommand{\N}{\mathbb{N}}
\newcommand{\Z}{\mathbb{Z}}
\newcommand{\Q}{\mathbb{Q}}
\newcommand{\R}{\mathbb{R}}
\newcommand{\C}{\mathbb{C}}
\newcommand{\F}{\mathbb{F}}
\newcommand{\A}{\mathbb{A}}
\newcommand{\PP}{\mathbb{P}}
\newcommand{\E}{\mathbb{E}}
\newcommand{\eps}{\varepsilon}
\newcommand{\dd}{\,\mathrm{d}}
\begin{document}
\section*{UnsolvedMath ID: 30004739}
\textbf{Source problem code:} catalogue entry
\subsection{Definitions and mathematical statement}
For a Euclidean billiard table \(D\), a billiard trajectory consists of straight segments with specular reflection at regular boundary points: the tangential component of velocity is preserved and the normal component is reversed. Write \(\operatorname{Per}_n(D)\) for trajectories returning to their initial position and direction after \(n\) reflections. This notation does not require \(n\) to be the least period; whether least period is intended is unresolved in the supplied record.

A one-parameter family of periodic trajectories is understood schematically as a nonconstant family \(\gamma_s\in\operatorname{Per}_n(D)\), with \(s\) varying over an interval. The topology, regularity, and identification of trajectories differing only by their starting point are not specified by the source. The trajectory-family parameter must be distinguished from the parameter describing deformation of the table.

\textbf{Question.} Can a billiard table be deformed nontrivially while simultaneously preserving a one-parameter family of \(2\)-periodic trajectories and a one-parameter family of \((2m+1)\)-periodic trajectories? Schematically, for a deformation \((D_t)_{t\in I}\), with \(I\) an interval containing \(0\), the requested persistence would give families
\[
\gamma_{2,t,s}\in\operatorname{Per}_2(D_t),\qquad
\gamma_{2m+1,t,u}\in\operatorname{Per}_{2m+1}(D_t)
\quad(t\in I),
\]
with \(s\) and \(u\) ranging over trajectory-family parameters, and with both families preserved throughout the deformation. This is a schematic formulation, not a complete quantified proposition: the source does not define preservation, admissible deformations, or nontriviality, and does not specify the range or quantification of \(m\).

\textbf{Supplied scope and qualifications.} The accompanying background discusses analytic-boundary deformations of a circle and a broader question for arbitrary convex domains. It does not fully specify the ambient dimension, boundary regularity in the broader class, or whether the intended problem concerns every eligible initial table, some initial table, or deformations from a circle. The relation between the stated periodic trajectory families and the rational caustics used in the supplied literature assessment also requires clarification. These missing choices prevent a unique precise formulation.

The supplied dated curation labels the problem partially solved and reports that, for deformations of a circle with analytic boundaries, simultaneous preservation of period-two and period-\((2m+1)\) rational caustics forces triviality, while the corresponding global question for arbitrary convex domains remains open. These are supplied literature claims, not independently verified mathematical conclusions; the exact hypotheses and meaning of triviality require review.
\subsection{Short English statement}
The question asks whether a deformation of a billiard table can preserve both a continuous family of two-periodic trajectories and a continuous family of odd-periodic trajectories. The admissible tables, deformations, meaning of preservation, and quantification of \(m\) remain unspecified. The supplied partial-solution assessment concerns analytic deformations from a circle and requires independent verification.
\subsection{Sources}
\begin{itemize}
\item Supplied local problem registry: permanent UnsolvedMath ID \(30004739\), source problem code \texttt{catalogue entry}, title \emph{Copreservation of Even and Odd Periodic Billiard Families}; collection: Oberwolfach Reports Open Problems. No individual proposer is supplied; the recorded proposed year is 2021.
\item Supplied report citation: \emph{Dynamische Systeme (2022)}, Oberwolfach Reports, DOI 10.4171/owr/2021/33, pp.\ 1735--1803. \url{https://doi.org/10.4171/owr/2021/33}. The record attributes the problem to workshop open-problem discussions documented in Oberwolfach Reports. The supplied citation year is retained without reconciling it with the proposed year or DOI year component.
\item Supplied literature reference labeled ``non-co-preservation theorem'': \url{https://arxiv.org/abs/2107.03499}. No authors, full title, or theorem number are supplied.
\item Supplied literature reference labeled ``period-two and period-three rigidity'': \url{https://arxiv.org/abs/1711.10541}. No authors, full title, or theorem number are supplied.
\item Supplied literature reference labeled ``Oberwolfach problem report'': \url{https://doi.org/10.4171/owr/2021/33}.
\item Literature assessment checked at \texttt{2026-08-22}: the supplied status is ``partially solved,'' with the analytic-circle rigidity claim and broader convex-domain open-status claim described above. This date and assessment are supplied dated curation derived from the OWR record, not independently verified source checks or a new literature assessment.
\end{itemize}
\end{document}
